\section*{Chapitre \ref{ch:fm:culture-and-decision-making} --- Culture et prise de décision}

\begin{solution}{exo:fm:culture-and-decision-making:1}
$(0.09+0.09+0.49)/3=0.223$~; une constante 0.5 obtient un score de 0.25.
\end{solution}

\begin{solution}{exo:fm:culture-and-decision-making:2}
0.70 pour les probabilités~; la moyenne des log-cotes donne 0.710, plus éloignée d'un demi.
\end{solution}

\begin{solution}{exo:fm:culture-and-decision-making:3}
Divisée par 5 si elles sont indépendantes~; avec $\rho=0.5$, la variance vaut $0.5+0.5/5=0.6$ de celle d'un seul prévisionniste, soit un facteur de seulement 1.67.
\end{solution}

\begin{solution}{exo:fm:culture-and-decision-making:4}
Son terme de fiabilité vaut 0.040 contre 0.017~: il étire ses log-cotes, si bien que ses prévisions extrêmes se réalisent beaucoup moins souvent qu'annoncé~; sa résolution
est à peu près la même.
\end{solution}

\begin{solution}{exo:fm:culture-and-decision-making:5}
Les prévisionnistes sont déjà trop confiants, si bien que pousser la moyenne plus loin d'un demi ajoute de la mauvaise calibration. L'extrémisation aide lorsque les prévisionnistes
disposent d'informations différentes et sont chacun bien calibrés, de sorte que leur moyenne est trop modérée.
\end{solution}

\begin{solution}{exo:fm:culture-and-decision-making:6}
Recueillir avant la réunion la probabilité que chacun attribue à la rupture de la stratégie, en discuter les éléments, recueillir de nouveau, puis enregistrer
l'agrégat des log-cotes dans le \omterm{def:fm:culture-and-decision-making:journal}{journal de décisions}, avec le critère d'abandon du chapitre~8.
\end{solution}

\begin{solution}{exo:fm:culture-and-decision-making:7}
Sur 362 paris d'espérance positive selon les probabilités propres du prévisionniste, 123 ont perdu~: une revue fondée sur les résultats qualifierait 34~\% d'entre eux d'erreurs.
\end{solution}

\begin{solution}{exo:fm:culture-and-decision-making:8}
Un résultat gagnant dit peu de chose de la décision~: il faut la juger d'après les probabilités et l'information disponibles à l'époque, consignées dans le \omterm{def:fm:culture-and-decision-making:journal}{journal de décisions}, et d'après la
calibration des prévisions passées de la personne.
\end{solution}

\begin{solution}{pb:fm:culture-and-decision-making:1}
\begin{enumerate}
\item La formation aux probabilités, le travail en équipe et le suivi ont amélioré la calibration et la résolution~; avec l'agrégation, les meilleures prévisions deux années de suite.
\item Voir la \cref{def:fm:culture-and-decision-making:journal}.
\item Voir la \cref{def:fm:culture-and-decision-making:brier}~; 0.25.
\item Fiabilité moins résolution plus incertitude~: la mauvaise calibration, le pouvoir discriminant entre événements et l'imprévisibilité propre des événements.
\item 0.240, 0.225, 0.252, 0.245 et 0.262~; la valeur vraie vaut 0.214.
\item Prévisionniste 1~: fiabilité 0.017, résolution 0.027~; prévisionniste 5~: 0.040 et 0.026~; la moyenne des log-cotes~: 0.010 et 0.039, avec une incertitude de 0.25.
\item Moyenne des probabilités 0.220, moyenne des log-cotes 0.223, moyenne extrémisée 0.237.
\item La courbe de calibration du prévisionniste trop confiant est plus plate que la diagonale~; celle de la moyenne s'en rapproche.
\item Voir la \cref{def:fm:culture-and-decision-making:agg}.
\item Voir la \cref{prop:fm:culture-and-decision-making:rho}.
\item 0.207 si indépendantes~; 0.224 à 0.5~; 0.237 à 1.
\item Les gens s'ancrent sur le premier chiffre, sur l'argument le plus bruyant ou sur l'avis de la personne la plus haut placée.
\item Des chiffres indépendants d'abord, la discussion ensuite, un second tour, puis agrégation des log-cotes.
\item Voir la \cref{def:fm:culture-and-decision-making:outcome}~; 34~\% des paris d'espérance positive ont perdu, et le prévisionniste 2 a remporté 66~\% de ses paris alors qu'il annonçait
77.7~\%.
\item Voir la \cref{def:fm:culture-and-decision-making:premortem}.
\item Par des scores sur de longues périodes et de nombreuses prévisions, non sur des résultats isolés (chapitre~10).
\item Le \omterm{def:fm:limits-and-capital-allocation:breach}{dépassement de limite}, les options, la probabilité que la position se rétablisse, le coût espéré de chaque option, qui a approuvé, et la date de revue.
\item La manipulation des scores par la prévision des seuls événements faciles, le repli vers un demi et la dissuasion de la contradiction~; noter toutes les décisions significatives, avec
la calibration et la résolution.
\item Dans le modèle du chapitre, agréger cinq prévisions indépendantes donne un score de 0.207, contre 0.224 pour une prévision discutée à $\rho=0.5$ (gain de
0.017) et 0.237 à $\rho=1$ (gain de 0.030).
\item Recueillir par écrit les probabilités de chacun avant toute discussion, agréger leurs log-cotes et noter tout le monde chaque trimestre~; conserver la
réunion, mais une fois les chiffres rassemblés.
\end{enumerate}
\end{solution}

\begin{solution}{iq:fm:culture-and-decision-making:1}
L'erreur quadratique moyenne des prévisions de probabilités sur des issues binaires~; 0.25 est le score obtenu en annonçant toujours un demi.
\iqlookfor{l'indice de référence sans compétence.}
\end{solution}

\begin{solution}{iq:fm:culture-and-decision-making:2}
On regarde les probabilités et l'information consignées au moment de la décision, ainsi que sa calibration sur de nombreuses transactions de ce type, et non ce résultat.
\iqlookfor{le processus plutôt que le résultat.}
\end{solution}

\begin{solution}{iq:fm:culture-and-decision-making:3}
Les log-cotes additionnent les éléments de manière symétrique et ne sont pas tirées vers un demi par un seul prévisionniste modéré~; moyenner les probabilités est trop
conservateur lorsque les prévisionnistes disposent d'informations différentes.
\iqlookfor{les éléments s'additionnent en log-cotes.}
\end{solution}

\begin{solution}{iq:fm:culture-and-decision-making:4}
La variance tombe à $0.5+0.5/5=0.6$ de celle d'une personne~: une réduction de 40~\%, contre 80~\% si les prévisions sont indépendantes.
\iqlookfor{$\rho+(1-\rho)/K$.}
\end{solution}

\begin{solution}{iq:fm:culture-and-decision-making:5}
Avant de s'engager, on demande à l'équipe de supposer que l'entrée a échoué au bout d'un an et d'écrire pourquoi~; on classe les raisons et on ajoute les plus probables aux
critères d'abandon et au chemin critique du plan.
\iqlookfor{échec imaginé, puis action.}
\end{solution}

\begin{solution}{iq:fm:culture-and-decision-making:6}
Regrouper les prévisions en classes~; développer $\frac1N\sum(p_i-y_i)^2$ autour de la fréquence observée de chaque classe et autour de la fréquence globale~; si les prévisions
sont constantes au sein de chaque classe, les termes croisés disparaissent, ce qui laisse fiabilité moins résolution plus incertitude.
\iqlookfor{le regroupement en classes et la disparition des termes croisés.}
\end{solution}
